% =========================================================
% ppt/src/minggu-14.tex
% Slide pertemuan 14
% Bab XV — Debugging pada Level Instruksi
%
% Sub-CPMK 4.1 dan 4.2 | CPMK04 | CPL03
% =========================================================

\documentclass[aspectratio=169,11pt]{beamer}
% =========================================================
% ppt/src/preamble-ppt.tex
% Tema, warna, paket, dan makro bersama seluruh deck slide.
%
% Dipakai oleh setiap berkas minggu-NN.tex dengan pola:
%   \documentclass[aspectratio=169,11pt]{beamer}
%   % =========================================================
% ppt/src/preamble-ppt.tex
% Tema, warna, paket, dan makro bersama seluruh deck slide.
%
% Dipakai oleh setiap berkas minggu-NN.tex dengan pola:
%   \documentclass[aspectratio=169,11pt]{beamer}
%   % =========================================================
% ppt/src/preamble-ppt.tex
% Tema, warna, paket, dan makro bersama seluruh deck slide.
%
% Dipakai oleh setiap berkas minggu-NN.tex dengan pola:
%   \documentclass[aspectratio=169,11pt]{beamer}
%   \input{preamble-ppt}
%
% Palet warna disalin dari buku_ajar/config/styles.tex agar
% identitas visual slide sama dengan buku ajar.
% =========================================================

% ---------------------------------------------------------
% 1. Paket
% ---------------------------------------------------------
\usepackage[T1]{fontenc}
\usepackage[utf8]{inputenc}
\usepackage[indonesian]{babel}
\usepackage{lmodern}
\usepackage{tikz}
\usetikzlibrary{arrows.meta, positioning, calc, shapes.geometric}
\usepackage{booktabs}
\usepackage{array}
\usepackage{amsmath}
\usepackage{amsfonts}
\usepackage{pifont}          % \ding untuk tanda checklist
\usepackage{listings}

% ---------------------------------------------------------
% 2. Palet warna OBE (identik dengan buku)
% ---------------------------------------------------------
\definecolor{obeBlue}{HTML}{1F3864}
\definecolor{obeBlueLight}{HTML}{DCE4F5}
\definecolor{obeTeal}{HTML}{0F6E6E}
\definecolor{obeTealLight}{HTML}{D9F0EE}
\definecolor{obeGreen}{HTML}{2E6B34}
\definecolor{obeGreenLight}{HTML}{DEEFDF}
\definecolor{obeAmber}{HTML}{8A5A00}
\definecolor{obeAmberLight}{HTML}{FDF1D6}
\definecolor{obeMaroon}{HTML}{7A2233}
\definecolor{obeMaroonLight}{HTML}{F7E1E5}
\definecolor{obeGray}{HTML}{4A4A4A}
\definecolor{obeGrayLight}{HTML}{F0F0F0}
\definecolor{codeBack}{HTML}{F7F7F4}
\definecolor{codeFrame}{HTML}{C9C9C0}

% ---------------------------------------------------------
% 3. Identitas mata kuliah (dipakai seluruh deck)
% ---------------------------------------------------------
\newcommand{\NamaMK}{Bahasa Pemrograman Tingkat Rendah}
\newcommand{\NamaMKShort}{Bhs.\ Pemrograman Tingkat Rendah}
\newcommand{\KodeMK}{MWCA010}
\newcommand{\Prodi}{Program Studi Teknik Informatika}
\newcommand{\Fakultas}{Fakultas Teknologi Informasi}
\newcommand{\PerguruanTinggi}{Universitas Bale Bandung}
\newcommand{\DosenPengampu}{Cecep Suwanda, S.Si., M.Kom.}
\newcommand{\SemesterTeks}{Semester Ganjil}

% \PertemuanTeks dan \JudulBab ditetapkan di masing-masing
% berkas minggu-NN.tex sebelum \begin{document}.
\providecommand{\PertemuanTeks}{Pertemuan}

% ---------------------------------------------------------
% 4. Tema beamer
% ---------------------------------------------------------
\usetheme{default}
\setbeamertemplate{navigation symbols}{}

\setbeamercolor{normal text}{fg=black!88}
\setbeamercolor{structure}{fg=obeBlue}
\setbeamercolor{alerted text}{fg=obeMaroon}

% --- judul slide: pita biru penuh lebar ---
\setbeamercolor{frametitle}{fg=white,bg=obeBlue}
\setbeamerfont{frametitle}{series=\bfseries,size=\large}
\setbeamertemplate{frametitle}{%
    \nointerlineskip%
    \begin{beamercolorbox}[%
        wd=\paperwidth,ht=2.9ex,dp=1.5ex,
        leftskip=0.4cm,rightskip=0.4cm]{frametitle}%
        \usebeamerfont{frametitle}\insertframetitle%
    \end{beamercolorbox}%
    \nointerlineskip%
    \hbox{\textcolor{obeTeal}{\rule{\paperwidth}{1.1pt}}}%
}

% --- kaki slide: identitas mata kuliah, pertemuan, nomor ---
\setbeamercolor{footline}{fg=white,bg=obeGray}
\setbeamerfont{footline}{size=\scriptsize,family=\sffamily}
\setbeamertemplate{footline}{%
    \leavevmode%
    \hbox{%
        \begin{beamercolorbox}[wd=0.42\paperwidth,ht=2.1ex,dp=1.1ex,
            leftskip=0.4cm]{footline}%
            \usebeamerfont{footline}%
            \KodeMK\ \textbullet\ \NamaMKShort%
        \end{beamercolorbox}%
        \begin{beamercolorbox}[wd=0.32\paperwidth,ht=2.1ex,dp=1.1ex,
            center]{footline}%
            \usebeamerfont{footline}\PertemuanTeks%
        \end{beamercolorbox}%
        \begin{beamercolorbox}[wd=0.26\paperwidth,ht=2.1ex,dp=1.1ex,
            rightskip=0.4cm]{footline}%
            \usebeamerfont{footline}%
            \insertframenumber\,/\,\inserttotalframenumber%
        \end{beamercolorbox}%
    }%
    \vskip0pt%
}

% --- daftar ---
\setbeamercolor{itemize item}{fg=obeTeal}
\setbeamercolor{itemize subitem}{fg=obeTeal}
\setbeamercolor{enumerate item}{fg=obeBlue}
\setbeamertemplate{itemize item}{\raise0.2ex\hbox{\tiny$\blacksquare$}}
\setbeamertemplate{itemize subitem}{\raise0.2ex\hbox{\scriptsize$\bullet$}}

% --- kotak (pengganti tcolorbox pada buku) ---
\setbeamercolor{block title}{bg=obeBlueLight,fg=obeBlue}
\setbeamercolor{block body}{bg=obeBlueLight!45,fg=black!88}
\setbeamercolor{block title alerted}{bg=obeMaroonLight,fg=obeMaroon}
\setbeamercolor{block body alerted}{bg=obeMaroonLight!45,fg=black!88}
\setbeamercolor{block title example}{bg=obeTealLight,fg=obeTeal}
\setbeamercolor{block body example}{bg=obeTealLight!45,fg=black!88}

% ---------------------------------------------------------
% 5. Makro notasi (sama persis dengan buku)
%
% Catatan: \note TIDAK didefinisikan di sini karena beamer
% sudah memakai \note untuk catatan pembicara. Untuk catatan
% di slide dipakai \catatan.
% ---------------------------------------------------------
\newcommand{\code}[1]{\texttt{#1}}
\newcommand{\asm}[1]{\texttt{#1}}
\newcommand{\reg}[1]{\texttt{#1}}
\newcommand{\flag}[1]{\texttt{#1}}
\newcommand{\dir}[1]{\texttt{#1}}
\newcommand{\file}[1]{\texttt{#1}}
\newcommand{\eng}[1]{\textit{#1}}

\newcommand{\hex}[1]{\texttt{#1h}}
\newcommand{\bin}[1]{\texttt{#1b}}
\newcommand{\dec}[1]{\texttt{#1}}

\newcommand{\important}[1]{\textbf{Penting:} #1}
\newcommand{\catatan}[1]{\textbf{Catatan:} #1}
\newcommand{\term}[1]{\textbf{#1}}

\newcommand{\CPL}[1]{\textbf{CPL#1}}
\newcommand{\CPMK}[1]{\textbf{CPMK#1}}
\newcommand{\SubCPMK}[1]{\textbf{Sub-CPMK #1}}

\newcommand{\checkbox}{\ding{113}\hspace{0.35em}}
\newcommand{\checked}{\ding{51}\hspace{0.35em}}
\newcommand{\crossed}{\ding{55}\hspace{0.35em}}

\newcommand{\flowarrow}{\ensuremath{\rightarrow}}

% ---------------------------------------------------------
% 6. Makro tata letak slide
% ---------------------------------------------------------

% Keterangan sumber materi (ketertelusuran ke buku)
\newcommand{\Sumber}[1]{%
    \vfill
    {\scriptsize\color{obeGray}Sumber: #1\par}%
}

% Frame judul deck: \JudulDeck{<judul bab>}{<pertemuan>}
\newcommand{\JudulDeck}[2]{%
\begin{frame}[plain]
    \vspace{0.25cm}
    {\color{obeBlue}\sffamily\bfseries\small \NamaMK\ (\KodeMK)\par}
    \vspace{0.12cm}
    {\color{obeGray}\sffamily\scriptsize
     \Prodi\ \textbullet\ \Fakultas\ \textbullet\ \PerguruanTinggi\par}
    \vfill
    {\color{obeBlue}\sffamily\LARGE\bfseries #1\par}
    \vspace{0.4cm}
    {\color{obeTeal}\sffamily\large #2\par}
    \vspace{0.45cm}
    {\color{obeTeal}\rule{\textwidth}{1pt}}
    \vfill
    {\color{obeGray}\sffamily\small Dosen Pengampu\par}
    \vspace{0.08cm}
    {\sffamily\bfseries \DosenPengampu\par}
    \vspace{0.2cm}
    {\color{obeGray}\sffamily\small \SemesterTeks\par}
    \vspace{0.25cm}
\end{frame}
}

% Frame pemisah bagian
\newcommand{\PemisahBagian}[1]{%
\begin{frame}[plain]
    \vfill
    \begin{center}
        {\color{obeTeal}\sffamily\small\bfseries BAGIAN\par}
        \vspace{0.35cm}
        {\color{obeBlue}\sffamily\Large\bfseries #1\par}
        \vspace{0.35cm}
        {\color{obeTeal}\rule{0.55\textwidth}{0.9pt}}
    \end{center}
    \vfill
\end{frame}
}

% ---------------------------------------------------------
% 7. Gaya penampil kode
% ---------------------------------------------------------
\input{listing-ppt}

%
% Palet warna disalin dari buku_ajar/config/styles.tex agar
% identitas visual slide sama dengan buku ajar.
% =========================================================

% ---------------------------------------------------------
% 1. Paket
% ---------------------------------------------------------
\usepackage[T1]{fontenc}
\usepackage[utf8]{inputenc}
\usepackage[indonesian]{babel}
\usepackage{lmodern}
\usepackage{tikz}
\usetikzlibrary{arrows.meta, positioning, calc, shapes.geometric}
\usepackage{booktabs}
\usepackage{array}
\usepackage{amsmath}
\usepackage{amsfonts}
\usepackage{pifont}          % \ding untuk tanda checklist
\usepackage{listings}

% ---------------------------------------------------------
% 2. Palet warna OBE (identik dengan buku)
% ---------------------------------------------------------
\definecolor{obeBlue}{HTML}{1F3864}
\definecolor{obeBlueLight}{HTML}{DCE4F5}
\definecolor{obeTeal}{HTML}{0F6E6E}
\definecolor{obeTealLight}{HTML}{D9F0EE}
\definecolor{obeGreen}{HTML}{2E6B34}
\definecolor{obeGreenLight}{HTML}{DEEFDF}
\definecolor{obeAmber}{HTML}{8A5A00}
\definecolor{obeAmberLight}{HTML}{FDF1D6}
\definecolor{obeMaroon}{HTML}{7A2233}
\definecolor{obeMaroonLight}{HTML}{F7E1E5}
\definecolor{obeGray}{HTML}{4A4A4A}
\definecolor{obeGrayLight}{HTML}{F0F0F0}
\definecolor{codeBack}{HTML}{F7F7F4}
\definecolor{codeFrame}{HTML}{C9C9C0}

% ---------------------------------------------------------
% 3. Identitas mata kuliah (dipakai seluruh deck)
% ---------------------------------------------------------
\newcommand{\NamaMK}{Bahasa Pemrograman Tingkat Rendah}
\newcommand{\NamaMKShort}{Bhs.\ Pemrograman Tingkat Rendah}
\newcommand{\KodeMK}{MWCA010}
\newcommand{\Prodi}{Program Studi Teknik Informatika}
\newcommand{\Fakultas}{Fakultas Teknologi Informasi}
\newcommand{\PerguruanTinggi}{Universitas Bale Bandung}
\newcommand{\DosenPengampu}{Cecep Suwanda, S.Si., M.Kom.}
\newcommand{\SemesterTeks}{Semester Ganjil}

% \PertemuanTeks dan \JudulBab ditetapkan di masing-masing
% berkas minggu-NN.tex sebelum \begin{document}.
\providecommand{\PertemuanTeks}{Pertemuan}

% ---------------------------------------------------------
% 4. Tema beamer
% ---------------------------------------------------------
\usetheme{default}
\setbeamertemplate{navigation symbols}{}

\setbeamercolor{normal text}{fg=black!88}
\setbeamercolor{structure}{fg=obeBlue}
\setbeamercolor{alerted text}{fg=obeMaroon}

% --- judul slide: pita biru penuh lebar ---
\setbeamercolor{frametitle}{fg=white,bg=obeBlue}
\setbeamerfont{frametitle}{series=\bfseries,size=\large}
\setbeamertemplate{frametitle}{%
    \nointerlineskip%
    \begin{beamercolorbox}[%
        wd=\paperwidth,ht=2.9ex,dp=1.5ex,
        leftskip=0.4cm,rightskip=0.4cm]{frametitle}%
        \usebeamerfont{frametitle}\insertframetitle%
    \end{beamercolorbox}%
    \nointerlineskip%
    \hbox{\textcolor{obeTeal}{\rule{\paperwidth}{1.1pt}}}%
}

% --- kaki slide: identitas mata kuliah, pertemuan, nomor ---
\setbeamercolor{footline}{fg=white,bg=obeGray}
\setbeamerfont{footline}{size=\scriptsize,family=\sffamily}
\setbeamertemplate{footline}{%
    \leavevmode%
    \hbox{%
        \begin{beamercolorbox}[wd=0.42\paperwidth,ht=2.1ex,dp=1.1ex,
            leftskip=0.4cm]{footline}%
            \usebeamerfont{footline}%
            \KodeMK\ \textbullet\ \NamaMKShort%
        \end{beamercolorbox}%
        \begin{beamercolorbox}[wd=0.32\paperwidth,ht=2.1ex,dp=1.1ex,
            center]{footline}%
            \usebeamerfont{footline}\PertemuanTeks%
        \end{beamercolorbox}%
        \begin{beamercolorbox}[wd=0.26\paperwidth,ht=2.1ex,dp=1.1ex,
            rightskip=0.4cm]{footline}%
            \usebeamerfont{footline}%
            \insertframenumber\,/\,\inserttotalframenumber%
        \end{beamercolorbox}%
    }%
    \vskip0pt%
}

% --- daftar ---
\setbeamercolor{itemize item}{fg=obeTeal}
\setbeamercolor{itemize subitem}{fg=obeTeal}
\setbeamercolor{enumerate item}{fg=obeBlue}
\setbeamertemplate{itemize item}{\raise0.2ex\hbox{\tiny$\blacksquare$}}
\setbeamertemplate{itemize subitem}{\raise0.2ex\hbox{\scriptsize$\bullet$}}

% --- kotak (pengganti tcolorbox pada buku) ---
\setbeamercolor{block title}{bg=obeBlueLight,fg=obeBlue}
\setbeamercolor{block body}{bg=obeBlueLight!45,fg=black!88}
\setbeamercolor{block title alerted}{bg=obeMaroonLight,fg=obeMaroon}
\setbeamercolor{block body alerted}{bg=obeMaroonLight!45,fg=black!88}
\setbeamercolor{block title example}{bg=obeTealLight,fg=obeTeal}
\setbeamercolor{block body example}{bg=obeTealLight!45,fg=black!88}

% ---------------------------------------------------------
% 5. Makro notasi (sama persis dengan buku)
%
% Catatan: \note TIDAK didefinisikan di sini karena beamer
% sudah memakai \note untuk catatan pembicara. Untuk catatan
% di slide dipakai \catatan.
% ---------------------------------------------------------
\newcommand{\code}[1]{\texttt{#1}}
\newcommand{\asm}[1]{\texttt{#1}}
\newcommand{\reg}[1]{\texttt{#1}}
\newcommand{\flag}[1]{\texttt{#1}}
\newcommand{\dir}[1]{\texttt{#1}}
\newcommand{\file}[1]{\texttt{#1}}
\newcommand{\eng}[1]{\textit{#1}}

\newcommand{\hex}[1]{\texttt{#1h}}
\newcommand{\bin}[1]{\texttt{#1b}}
\newcommand{\dec}[1]{\texttt{#1}}

\newcommand{\important}[1]{\textbf{Penting:} #1}
\newcommand{\catatan}[1]{\textbf{Catatan:} #1}
\newcommand{\term}[1]{\textbf{#1}}

\newcommand{\CPL}[1]{\textbf{CPL#1}}
\newcommand{\CPMK}[1]{\textbf{CPMK#1}}
\newcommand{\SubCPMK}[1]{\textbf{Sub-CPMK #1}}

\newcommand{\checkbox}{\ding{113}\hspace{0.35em}}
\newcommand{\checked}{\ding{51}\hspace{0.35em}}
\newcommand{\crossed}{\ding{55}\hspace{0.35em}}

\newcommand{\flowarrow}{\ensuremath{\rightarrow}}

% ---------------------------------------------------------
% 6. Makro tata letak slide
% ---------------------------------------------------------

% Keterangan sumber materi (ketertelusuran ke buku)
\newcommand{\Sumber}[1]{%
    \vfill
    {\scriptsize\color{obeGray}Sumber: #1\par}%
}

% Frame judul deck: \JudulDeck{<judul bab>}{<pertemuan>}
\newcommand{\JudulDeck}[2]{%
\begin{frame}[plain]
    \vspace{0.25cm}
    {\color{obeBlue}\sffamily\bfseries\small \NamaMK\ (\KodeMK)\par}
    \vspace{0.12cm}
    {\color{obeGray}\sffamily\scriptsize
     \Prodi\ \textbullet\ \Fakultas\ \textbullet\ \PerguruanTinggi\par}
    \vfill
    {\color{obeBlue}\sffamily\LARGE\bfseries #1\par}
    \vspace{0.4cm}
    {\color{obeTeal}\sffamily\large #2\par}
    \vspace{0.45cm}
    {\color{obeTeal}\rule{\textwidth}{1pt}}
    \vfill
    {\color{obeGray}\sffamily\small Dosen Pengampu\par}
    \vspace{0.08cm}
    {\sffamily\bfseries \DosenPengampu\par}
    \vspace{0.2cm}
    {\color{obeGray}\sffamily\small \SemesterTeks\par}
    \vspace{0.25cm}
\end{frame}
}

% Frame pemisah bagian
\newcommand{\PemisahBagian}[1]{%
\begin{frame}[plain]
    \vfill
    \begin{center}
        {\color{obeTeal}\sffamily\small\bfseries BAGIAN\par}
        \vspace{0.35cm}
        {\color{obeBlue}\sffamily\Large\bfseries #1\par}
        \vspace{0.35cm}
        {\color{obeTeal}\rule{0.55\textwidth}{0.9pt}}
    \end{center}
    \vfill
\end{frame}
}

% ---------------------------------------------------------
% 7. Gaya penampil kode
% ---------------------------------------------------------
% =========================================================
% ppt/src/listing-ppt.tex
% Gaya penampil kode untuk slide.
%
% Diturunkan dari buku_ajar/config/listings.tex, tetapi
% disesuaikan untuk slide: ukuran lebih kecil (karena frame
% 16:9 dan kode harus utuh terlihat) dan jarak lebih rapat.
% =========================================================

% ---------------------------------------------------------
% 1. Definisi bahasa MASM/TASM 8086 (salinan dari buku)
% ---------------------------------------------------------
\lstdefinelanguage{MASM}{
    morekeywords={
        MOV, XCHG, LEA, LDS, LES, PUSH, POP, PUSHA, POPA, PUSHF, POPF,
        LAHF, SAHF, XLAT, IN, OUT, CBW, CWD,
        ADD, ADC, SUB, SBB, INC, DEC, NEG, CMP, MUL, IMUL, DIV, IDIV,
        AAA, AAS, AAM, AAD, DAA, DAS,
        AND, OR, XOR, NOT, TEST, SHL, SHR, SAL, SAR, ROL, ROR, RCL, RCR,
        JMP, JE, JZ, JNE, JNZ, JG, JGE, JL, JLE, JA, JAE, JB, JBE,
        JC, JNC, JS, JNS, JO, JNO, JP, JPE, JNP, JPO, JCXZ,
        LOOP, LOOPE, LOOPZ, LOOPNE, LOOPNZ, CALL, RET, INT, IRET, INTO,
        HLT, NOP, CLI, STI, CLD, STD, CLC, STC, CMC, WAIT, ESC, LOCK,
        DB, DW, DD, DQ, DT, EQU, ORG, SEGMENT, ENDS, ASSUME, END,
        PROC, ENDP, MACRO, ENDM, LOCAL, PUBLIC, EXTRN, EXTERN,
        INCLUDE, INCLUDELIB, MODEL, STACK, CODE, DATA, CONST, ALIGN, EVEN,
        OFFSET, PTR, BYTE, WORD, DWORD, QWORD, TBYTE,
        NEAR, FAR, SHORT, LABEL, TYPE, SIZE, LENGTH, DUP, THIS,
        AX, BX, CX, DX, AL, AH, BL, BH, CL, CH, DL, DH,
        SP, BP, SI, DI, IP, CS, DS, SS, ES, FLAGS,
        CF, PF, AF, ZF, TF, IF, DF, OF
    },
    sensitive=false,
    morecomment=[l]{;},
    morestring=[b]',
    morestring=[b]"
}

% ---------------------------------------------------------
% 2. Gaya kode assembly untuk slide
% ---------------------------------------------------------
\lstdefinestyle{asmslide}{
    language=MASM,
    basicstyle=\ttfamily\scriptsize,
    keywordstyle=\color{obeBlue}\bfseries,
    commentstyle=\color{obeGreen}\itshape,
    stringstyle=\color{obeMaroon},
    numberstyle=\tiny\color{obeGray},
    numbers=left,
    numbersep=5pt,
    stepnumber=1,
    backgroundcolor=\color{codeBack},
    frame=single,
    rulecolor=\color{codeFrame},
    framesep=4pt,
    breaklines=true,
    breakatwhitespace=false,
    % CATATAN PENTING: opsi `postbreak` TIDAK boleh dipakai di sini.
    % Di beamer, `breaklines=true` bersama `postbreak=<kode>` membuat
    % listings gagal dirakit dengan galat
    %   "Missing number, treated as zero <to be read again> \kern".
    % Karena \lstset{style=asmslide} bersifat global, opsi postbreak
    % juga bocor ke gaya lain (mis. outslide/keluaran) dan merusaknya.
    % Di buku ajar (kelas `book`) justru sebaliknya: postbreak aman dan
    % dipakai. Jadi larangan ini khusus untuk slide.
    showstringspaces=false,
    tabsize=4,
    columns=fullflexible,
    keepspaces=true,
    upquote=true,
    xleftmargin=12pt,
    aboveskip=0.4em,
    belowskip=0.4em
}

% ---------------------------------------------------------
% 3. Gaya keluaran program / sesi debugger untuk slide
% ---------------------------------------------------------
\lstdefinestyle{outslide}{
    basicstyle=\ttfamily\scriptsize,
    keywordstyle=\color{black},
    commentstyle=\color{obeGray},
    stringstyle=\color{black},
    numbers=none,
    backgroundcolor=\color{obeGrayLight},
    frame=single,
    rulecolor=\color{codeFrame},
    framesep=4pt,
    breaklines=true,
    showstringspaces=false,
    columns=fullflexible,
    keepspaces=true,
    upquote=true,
    xleftmargin=4pt,
    aboveskip=0.4em,
    belowskip=0.4em
}

\lstset{style=asmslide}

% ---------------------------------------------------------
% 4. Environment ringkas untuk dipakai di dalam frame
%
% Frame yang memuat kedua environment ini WAJIB memakai
% opsi [fragile], karena isinya verbatim.
% ---------------------------------------------------------
\lstnewenvironment{kodeasm}[1][]{%
    \lstset{style=asmslide, #1}%
}{}

\lstnewenvironment{keluaran}[1][]{%
    \lstset{style=outslide, #1}%
}{}


%
% Palet warna disalin dari buku_ajar/config/styles.tex agar
% identitas visual slide sama dengan buku ajar.
% =========================================================

% ---------------------------------------------------------
% 1. Paket
% ---------------------------------------------------------
\usepackage[T1]{fontenc}
\usepackage[utf8]{inputenc}
\usepackage[indonesian]{babel}
\usepackage{lmodern}
\usepackage{tikz}
\usetikzlibrary{arrows.meta, positioning, calc, shapes.geometric}
\usepackage{booktabs}
\usepackage{array}
\usepackage{amsmath}
\usepackage{amsfonts}
\usepackage{pifont}          % \ding untuk tanda checklist
\usepackage{listings}

% ---------------------------------------------------------
% 2. Palet warna OBE (identik dengan buku)
% ---------------------------------------------------------
\definecolor{obeBlue}{HTML}{1F3864}
\definecolor{obeBlueLight}{HTML}{DCE4F5}
\definecolor{obeTeal}{HTML}{0F6E6E}
\definecolor{obeTealLight}{HTML}{D9F0EE}
\definecolor{obeGreen}{HTML}{2E6B34}
\definecolor{obeGreenLight}{HTML}{DEEFDF}
\definecolor{obeAmber}{HTML}{8A5A00}
\definecolor{obeAmberLight}{HTML}{FDF1D6}
\definecolor{obeMaroon}{HTML}{7A2233}
\definecolor{obeMaroonLight}{HTML}{F7E1E5}
\definecolor{obeGray}{HTML}{4A4A4A}
\definecolor{obeGrayLight}{HTML}{F0F0F0}
\definecolor{codeBack}{HTML}{F7F7F4}
\definecolor{codeFrame}{HTML}{C9C9C0}

% ---------------------------------------------------------
% 3. Identitas mata kuliah (dipakai seluruh deck)
% ---------------------------------------------------------
\newcommand{\NamaMK}{Bahasa Pemrograman Tingkat Rendah}
\newcommand{\NamaMKShort}{Bhs.\ Pemrograman Tingkat Rendah}
\newcommand{\KodeMK}{MWCA010}
\newcommand{\Prodi}{Program Studi Teknik Informatika}
\newcommand{\Fakultas}{Fakultas Teknologi Informasi}
\newcommand{\PerguruanTinggi}{Universitas Bale Bandung}
\newcommand{\DosenPengampu}{Cecep Suwanda, S.Si., M.Kom.}
\newcommand{\SemesterTeks}{Semester Ganjil}

% \PertemuanTeks dan \JudulBab ditetapkan di masing-masing
% berkas minggu-NN.tex sebelum \begin{document}.
\providecommand{\PertemuanTeks}{Pertemuan}

% ---------------------------------------------------------
% 4. Tema beamer
% ---------------------------------------------------------
\usetheme{default}
\setbeamertemplate{navigation symbols}{}

\setbeamercolor{normal text}{fg=black!88}
\setbeamercolor{structure}{fg=obeBlue}
\setbeamercolor{alerted text}{fg=obeMaroon}

% --- judul slide: pita biru penuh lebar ---
\setbeamercolor{frametitle}{fg=white,bg=obeBlue}
\setbeamerfont{frametitle}{series=\bfseries,size=\large}
\setbeamertemplate{frametitle}{%
    \nointerlineskip%
    \begin{beamercolorbox}[%
        wd=\paperwidth,ht=2.9ex,dp=1.5ex,
        leftskip=0.4cm,rightskip=0.4cm]{frametitle}%
        \usebeamerfont{frametitle}\insertframetitle%
    \end{beamercolorbox}%
    \nointerlineskip%
    \hbox{\textcolor{obeTeal}{\rule{\paperwidth}{1.1pt}}}%
}

% --- kaki slide: identitas mata kuliah, pertemuan, nomor ---
\setbeamercolor{footline}{fg=white,bg=obeGray}
\setbeamerfont{footline}{size=\scriptsize,family=\sffamily}
\setbeamertemplate{footline}{%
    \leavevmode%
    \hbox{%
        \begin{beamercolorbox}[wd=0.42\paperwidth,ht=2.1ex,dp=1.1ex,
            leftskip=0.4cm]{footline}%
            \usebeamerfont{footline}%
            \KodeMK\ \textbullet\ \NamaMKShort%
        \end{beamercolorbox}%
        \begin{beamercolorbox}[wd=0.32\paperwidth,ht=2.1ex,dp=1.1ex,
            center]{footline}%
            \usebeamerfont{footline}\PertemuanTeks%
        \end{beamercolorbox}%
        \begin{beamercolorbox}[wd=0.26\paperwidth,ht=2.1ex,dp=1.1ex,
            rightskip=0.4cm]{footline}%
            \usebeamerfont{footline}%
            \insertframenumber\,/\,\inserttotalframenumber%
        \end{beamercolorbox}%
    }%
    \vskip0pt%
}

% --- daftar ---
\setbeamercolor{itemize item}{fg=obeTeal}
\setbeamercolor{itemize subitem}{fg=obeTeal}
\setbeamercolor{enumerate item}{fg=obeBlue}
\setbeamertemplate{itemize item}{\raise0.2ex\hbox{\tiny$\blacksquare$}}
\setbeamertemplate{itemize subitem}{\raise0.2ex\hbox{\scriptsize$\bullet$}}

% --- kotak (pengganti tcolorbox pada buku) ---
\setbeamercolor{block title}{bg=obeBlueLight,fg=obeBlue}
\setbeamercolor{block body}{bg=obeBlueLight!45,fg=black!88}
\setbeamercolor{block title alerted}{bg=obeMaroonLight,fg=obeMaroon}
\setbeamercolor{block body alerted}{bg=obeMaroonLight!45,fg=black!88}
\setbeamercolor{block title example}{bg=obeTealLight,fg=obeTeal}
\setbeamercolor{block body example}{bg=obeTealLight!45,fg=black!88}

% ---------------------------------------------------------
% 5. Makro notasi (sama persis dengan buku)
%
% Catatan: \note TIDAK didefinisikan di sini karena beamer
% sudah memakai \note untuk catatan pembicara. Untuk catatan
% di slide dipakai \catatan.
% ---------------------------------------------------------
\newcommand{\code}[1]{\texttt{#1}}
\newcommand{\asm}[1]{\texttt{#1}}
\newcommand{\reg}[1]{\texttt{#1}}
\newcommand{\flag}[1]{\texttt{#1}}
\newcommand{\dir}[1]{\texttt{#1}}
\newcommand{\file}[1]{\texttt{#1}}
\newcommand{\eng}[1]{\textit{#1}}

\newcommand{\hex}[1]{\texttt{#1h}}
\newcommand{\bin}[1]{\texttt{#1b}}
\newcommand{\dec}[1]{\texttt{#1}}

\newcommand{\important}[1]{\textbf{Penting:} #1}
\newcommand{\catatan}[1]{\textbf{Catatan:} #1}
\newcommand{\term}[1]{\textbf{#1}}

\newcommand{\CPL}[1]{\textbf{CPL#1}}
\newcommand{\CPMK}[1]{\textbf{CPMK#1}}
\newcommand{\SubCPMK}[1]{\textbf{Sub-CPMK #1}}

\newcommand{\checkbox}{\ding{113}\hspace{0.35em}}
\newcommand{\checked}{\ding{51}\hspace{0.35em}}
\newcommand{\crossed}{\ding{55}\hspace{0.35em}}

\newcommand{\flowarrow}{\ensuremath{\rightarrow}}

% ---------------------------------------------------------
% 6. Makro tata letak slide
% ---------------------------------------------------------

% Keterangan sumber materi (ketertelusuran ke buku)
\newcommand{\Sumber}[1]{%
    \vfill
    {\scriptsize\color{obeGray}Sumber: #1\par}%
}

% Frame judul deck: \JudulDeck{<judul bab>}{<pertemuan>}
\newcommand{\JudulDeck}[2]{%
\begin{frame}[plain]
    \vspace{0.25cm}
    {\color{obeBlue}\sffamily\bfseries\small \NamaMK\ (\KodeMK)\par}
    \vspace{0.12cm}
    {\color{obeGray}\sffamily\scriptsize
     \Prodi\ \textbullet\ \Fakultas\ \textbullet\ \PerguruanTinggi\par}
    \vfill
    {\color{obeBlue}\sffamily\LARGE\bfseries #1\par}
    \vspace{0.4cm}
    {\color{obeTeal}\sffamily\large #2\par}
    \vspace{0.45cm}
    {\color{obeTeal}\rule{\textwidth}{1pt}}
    \vfill
    {\color{obeGray}\sffamily\small Dosen Pengampu\par}
    \vspace{0.08cm}
    {\sffamily\bfseries \DosenPengampu\par}
    \vspace{0.2cm}
    {\color{obeGray}\sffamily\small \SemesterTeks\par}
    \vspace{0.25cm}
\end{frame}
}

% Frame pemisah bagian
\newcommand{\PemisahBagian}[1]{%
\begin{frame}[plain]
    \vfill
    \begin{center}
        {\color{obeTeal}\sffamily\small\bfseries BAGIAN\par}
        \vspace{0.35cm}
        {\color{obeBlue}\sffamily\Large\bfseries #1\par}
        \vspace{0.35cm}
        {\color{obeTeal}\rule{0.55\textwidth}{0.9pt}}
    \end{center}
    \vfill
\end{frame}
}

% ---------------------------------------------------------
% 7. Gaya penampil kode
% ---------------------------------------------------------
% =========================================================
% ppt/src/listing-ppt.tex
% Gaya penampil kode untuk slide.
%
% Diturunkan dari buku_ajar/config/listings.tex, tetapi
% disesuaikan untuk slide: ukuran lebih kecil (karena frame
% 16:9 dan kode harus utuh terlihat) dan jarak lebih rapat.
% =========================================================

% ---------------------------------------------------------
% 1. Definisi bahasa MASM/TASM 8086 (salinan dari buku)
% ---------------------------------------------------------
\lstdefinelanguage{MASM}{
    morekeywords={
        MOV, XCHG, LEA, LDS, LES, PUSH, POP, PUSHA, POPA, PUSHF, POPF,
        LAHF, SAHF, XLAT, IN, OUT, CBW, CWD,
        ADD, ADC, SUB, SBB, INC, DEC, NEG, CMP, MUL, IMUL, DIV, IDIV,
        AAA, AAS, AAM, AAD, DAA, DAS,
        AND, OR, XOR, NOT, TEST, SHL, SHR, SAL, SAR, ROL, ROR, RCL, RCR,
        JMP, JE, JZ, JNE, JNZ, JG, JGE, JL, JLE, JA, JAE, JB, JBE,
        JC, JNC, JS, JNS, JO, JNO, JP, JPE, JNP, JPO, JCXZ,
        LOOP, LOOPE, LOOPZ, LOOPNE, LOOPNZ, CALL, RET, INT, IRET, INTO,
        HLT, NOP, CLI, STI, CLD, STD, CLC, STC, CMC, WAIT, ESC, LOCK,
        DB, DW, DD, DQ, DT, EQU, ORG, SEGMENT, ENDS, ASSUME, END,
        PROC, ENDP, MACRO, ENDM, LOCAL, PUBLIC, EXTRN, EXTERN,
        INCLUDE, INCLUDELIB, MODEL, STACK, CODE, DATA, CONST, ALIGN, EVEN,
        OFFSET, PTR, BYTE, WORD, DWORD, QWORD, TBYTE,
        NEAR, FAR, SHORT, LABEL, TYPE, SIZE, LENGTH, DUP, THIS,
        AX, BX, CX, DX, AL, AH, BL, BH, CL, CH, DL, DH,
        SP, BP, SI, DI, IP, CS, DS, SS, ES, FLAGS,
        CF, PF, AF, ZF, TF, IF, DF, OF
    },
    sensitive=false,
    morecomment=[l]{;},
    morestring=[b]',
    morestring=[b]"
}

% ---------------------------------------------------------
% 2. Gaya kode assembly untuk slide
% ---------------------------------------------------------
\lstdefinestyle{asmslide}{
    language=MASM,
    basicstyle=\ttfamily\scriptsize,
    keywordstyle=\color{obeBlue}\bfseries,
    commentstyle=\color{obeGreen}\itshape,
    stringstyle=\color{obeMaroon},
    numberstyle=\tiny\color{obeGray},
    numbers=left,
    numbersep=5pt,
    stepnumber=1,
    backgroundcolor=\color{codeBack},
    frame=single,
    rulecolor=\color{codeFrame},
    framesep=4pt,
    breaklines=true,
    breakatwhitespace=false,
    % CATATAN PENTING: opsi `postbreak` TIDAK boleh dipakai di sini.
    % Di beamer, `breaklines=true` bersama `postbreak=<kode>` membuat
    % listings gagal dirakit dengan galat
    %   "Missing number, treated as zero <to be read again> \kern".
    % Karena \lstset{style=asmslide} bersifat global, opsi postbreak
    % juga bocor ke gaya lain (mis. outslide/keluaran) dan merusaknya.
    % Di buku ajar (kelas `book`) justru sebaliknya: postbreak aman dan
    % dipakai. Jadi larangan ini khusus untuk slide.
    showstringspaces=false,
    tabsize=4,
    columns=fullflexible,
    keepspaces=true,
    upquote=true,
    xleftmargin=12pt,
    aboveskip=0.4em,
    belowskip=0.4em
}

% ---------------------------------------------------------
% 3. Gaya keluaran program / sesi debugger untuk slide
% ---------------------------------------------------------
\lstdefinestyle{outslide}{
    basicstyle=\ttfamily\scriptsize,
    keywordstyle=\color{black},
    commentstyle=\color{obeGray},
    stringstyle=\color{black},
    numbers=none,
    backgroundcolor=\color{obeGrayLight},
    frame=single,
    rulecolor=\color{codeFrame},
    framesep=4pt,
    breaklines=true,
    showstringspaces=false,
    columns=fullflexible,
    keepspaces=true,
    upquote=true,
    xleftmargin=4pt,
    aboveskip=0.4em,
    belowskip=0.4em
}

\lstset{style=asmslide}

% ---------------------------------------------------------
% 4. Environment ringkas untuk dipakai di dalam frame
%
% Frame yang memuat kedua environment ini WAJIB memakai
% opsi [fragile], karena isinya verbatim.
% ---------------------------------------------------------
\lstnewenvironment{kodeasm}[1][]{%
    \lstset{style=asmslide, #1}%
}{}

\lstnewenvironment{keluaran}[1][]{%
    \lstset{style=outslide, #1}%
}{}



\renewcommand{\PertemuanTeks}{Pertemuan 14}

\begin{document}

% =========================================================
% 1. JUDUL
% =========================================================
\JudulDeck{Debugging pada Level Instruksi}{Pertemuan 14}

% =========================================================
% 2. CAPAIAN DAN TUJUAN
% =========================================================
\begin{frame}{Capaian dan Tujuan Pembelajaran}
\begin{block}{Sub-CPMK 4.1 dan 4.2}
\scriptsize
\SubCPMK{4.1} --- Melakukan penelusuran program pada level instruksi
(\eng{single-step}, \eng{trace}, \eng{breakpoint}) dan memantau
perubahan register serta \eng{flag}.

\vspace{0.3em}
\SubCPMK{4.2} --- Memeriksa dan menganalisis isi memori dan
\eng{stack} (\eng{dump}) untuk menemukan kesalahan logika, kesalahan
segmentasi, dan \eng{overflow}.
\end{block}

{\small\textbf{CPMK yang didukung:} \CPMK{04}\quad\textbullet\quad
\textbf{CPL yang didukung:} \CPL{03}\par}

\begin{exampleblock}{Setelah pertemuan ini, mahasiswa mampu:}
\scriptsize
\begin{enumerate}
    \item membedakan lima jenis kesalahan program: sintaks, saat
          berjalan, logika, segmentasi, dan \eng{overflow};
    \item memilih alat \eng{debugging} yang sesuai: MS-DEBUG,
          \eng{debugger} GUI Turbo Assembler, atau Emu8086;
    \item memakai perintah \code{R}, \code{U}, \code{T}, \code{P},
          \code{G}, \code{D}, \code{E}, \code{A}, \code{N},
          \code{L}, dan \code{W} pada DEBUG;
    \item menjelaskan perbedaan \code{T} dan \code{P} ketika
          menjumpai \asm{CALL} dan \code{INT};
    \item mengisi tabel \eng{trace} register dan \eng{flag} setiap
          langkah eksekusi;
    \item membaca keluaran \code{D} serta menafsirkan bingkai
          \eng{stack};
    \item menerapkan metodologi \eng{debugging} sistematis untuk
          menemukan tiga bug khas.
\end{enumerate}
\end{exampleblock}
\end{frame}

% =========================================================
% 3. PETA MATERI
% =========================================================
\begin{frame}{Peta Materi}
\small
\begin{enumerate}
    \item \textbf{Taksonomi Kesalahan Program dan Alat Debugging}\\
          {\footnotesize lima jenis kesalahan; gejala dan cara
          penemuannya; MS-DEBUG, GUI Turbo Assembler, dan Emu8086}

    \vspace{0.25em}
    \item \textbf{Menelusuri Program dengan Perintah DEBUG}\\
          {\footnotesize \code{N}, \code{L}, \code{W}, \code{R},
          \code{U}, \code{T}, \code{P}, \code{G}, \code{D},
          \code{E}, \code{A}; perbedaan \code{T} dan \code{P};
          \eng{breakpoint}; tabel \eng{trace}}

    \vspace{0.25em}
    \item \textbf{Analisis Memori dan Stack serta Metodologi
          Debugging}\\
          {\footnotesize membaca \eng{dump}; bingkai \eng{stack};
          deteksi \eng{overflow}; enam tahap metodologi; tiga studi
          kasus}
\end{enumerate}

\vspace{0.35em}
\begin{block}{Praktik pada pertemuan ini}
\small Menelusuri \file{bug-loop.asm} dan \file{bug-stack.asm} dengan
MS-DEBUG, mengisi tabel \eng{trace}, lalu memperbaiki kedua program.
\end{block}
\end{frame}

% =========================================================
% BAGIAN 1
% =========================================================
\PemisahBagian{Taksonomi Kesalahan Program dan Alat Debugging}

% ---------------------------------------------------------
\begin{frame}{Debugging sebagai Proses Sistematis}
\small
Menulis program assembly berarti berhadapan langsung dengan register,
\eng{flag}, dan alamat memori tanpa lapisan pelindung apa pun.

\vspace{0.3em}
\begin{exampleblock}{Definisi}
\small \term{Debugging} adalah proses sistematis untuk menemukan
penyebab suatu kesalahan (\eng{bug}) pada program, membuktikan
penyebab itu melalui bukti pengamatan, lalu memperbaiki dan
memverifikasi hasil perbaikan.
\end{exampleblock}

\vspace{0.25em}
\small
\begin{itemize}
    \item \eng{Debugging} bukan kegiatan menebak-nebak, melainkan
          kegiatan mengumpulkan bukti pada level instruksi.
    \item Bekal utamanya adalah tiga pengetahuan dari bab
          sebelumnya: cara kerja register, pengambilan keputusan
          berdasarkan \eng{flag}, dan pemakaian \eng{stack} oleh
          prosedur.
\end{itemize}
\Sumber{Bab XV, Taksonomi Kesalahan Program dan Alat Debugging}
\end{frame}

% ---------------------------------------------------------
\begin{frame}{Lima Jenis Kesalahan Program}
\footnotesize
\begin{enumerate}
    \item \term{Kesalahan sintaks} --- penulisan yang membuat
          assembler menolak berkas sumber, sehingga program
          \emph{tidak} dihasilkan (mis. \code{MOVV} alih-alih
          \asm{MOV}).

    \vspace{0.2em}
    \item \term{Kesalahan saat berjalan} --- baru muncul ketika
          program dijalankan, misalnya pembagian dengan nol.

    \vspace{0.2em}
    \item \term{Kesalahan logika} --- program berjalan sampai selesai
          tetapi hasilnya salah; assembler dan sistem operasi
          sama-sama diam.

    \vspace{0.2em}
    \item \term{Kesalahan segmentasi} --- program mengakses alamat di
          luar segmen yang seharusnya, misalnya \reg{DS} tidak pernah
          diarahkan ke segmen data.

    \vspace{0.2em}
    \item \term{Overflow} --- hasil perhitungan atau pemakaian
          \eng{stack} melebihi daya tampungnya, misalnya melebihi
          \hex{FFFF} atau jumlah \asm{PUSH} tidak seimbang dengan
          \asm{POP}.
\end{enumerate}
\Sumber{Bab XV, Lima Jenis Kesalahan Program}
\end{frame}

% ---------------------------------------------------------
\begin{frame}{Gejala dan Cara Menemukan Tiap Jenis}
\scriptsize
\begin{center}
\begin{tabular}{@{}p{1.9cm}p{3.3cm}p{2.9cm}p{3.5cm}@{}}
    \toprule
    \textbf{Jenis} & \textbf{Gejala} & \textbf{Kapan terdeteksi} &
    \textbf{Cara menemukan} \\
    \midrule
    Sintaks & Assembler menolak berkas; tidak ada \file{.obj} &
    Sebelum program dijalankan & Baca pesan assembler beserta nomor
    barisnya \\
    \addlinespace[2pt]
    Saat berjalan & Program berhenti di tengah jalan; pesan sistem
    operasi & Ketika dijalankan & Jalankan bertahap sampai titik
    berhenti \\
    \addlinespace[2pt]
    Logika & Program selesai normal, hasil salah & Tidak terdeteksi
    otomatis & Tabel \eng{trace} register dan \eng{flag} per langkah \\
    \addlinespace[2pt]
    Segmentasi & Karakter sampah, nilai berubah sendiri, program
    macet & Ketika dijalankan & Periksa \reg{DS}, \reg{ES},
    \reg{SS}, dan batas indeks melalui \code{D} \\
    \addlinespace[2pt]
    \eng{Overflow} & Hasil aritmatika terpotong; \asm{RET} melompat
    ke alamat aneh & Ketika dijalankan & Periksa \flag{CF},
    \flag{OF}, dan pasangan \eng{dump} \eng{stack} \\
    \bottomrule
\end{tabular}
\end{center}

\vspace{0.25em}
{\scriptsize\color{obeGray}\important{Hanya kesalahan sintaks yang
dapat dipercayakan kepada assembler.} Empat jenis lainnya menuntut
pengamatan langsung terhadap keadaan prosesor saat program
berjalan.\par}
\Sumber{Bab XV, Tabel taksonomi kesalahan}
\end{frame}

% ---------------------------------------------------------
\begin{frame}{Tiga Alat Debugging yang Tersedia}
\small
Ketiganya mampu melakukan hal yang sama: menghentikan program,
memeriksa register, dan memeriksa memori. Perbedaannya terletak pada
kemudahan pemakaian dan kedekatan dengan mesin sesungguhnya.

\vspace{0.3em}
\begin{enumerate}
    \item \term{MS-DEBUG di DOSBox.} Program bawaan DOS yang seluruh
          interaksinya berupa baris perintah. Yang terlihat
          benar-benar isi memori, bukan tampilan yang disederhanakan.

    \vspace{0.2em}
    \item \term{\eng{Debugger} terpadu GUI Turbo Assembler.} Satu
          jendela terpadu memuat kode sumber, register, dan memori
          sekaligus; program dijalankan baris per baris dengan
          menekan tombol.

    \vspace{0.2em}
    \item \term{Jendela \eng{debug} Emu8086.} Register, \eng{flag},
          memori, dan \eng{stack} tampil sebagai panel terpisah yang
          selalu terlihat sehingga cocok untuk pemula.
\end{enumerate}

\vspace{0.15em}
\begin{block}{Mengapa MS-DEBUG ditekankan}
\small DEBUG memperlihatkan keadaan mesin apa adanya, dan kemampuan
membaca keluarannya berlaku umum bagi kedua alat grafis lainnya.
\end{block}
\Sumber{Bab XV, Alat Debugging yang Tersedia}
\end{frame}

% ---------------------------------------------------------
\begin{frame}{Perbandingan Kemampuan Tiga Alat}
\scriptsize
\begin{center}
\begin{tabular}{@{}p{2.3cm}p{3.4cm}p{3.6cm}p{3.0cm}@{}}
    \toprule
    \textbf{Kemampuan} & \textbf{MS-DEBUG (DOSBox)} &
    \textbf{GUI Turbo Assembler} & \textbf{Emu8086} \\
    \midrule
    Tampilan & Teks, baris perintah & Grafis, jendela terpadu &
    Grafis, panel terpisah \\
    \addlinespace[2pt]
    Langkah per instruksi & \code{T} dan \code{P} & Tombol
    \file{Step} & Tombol \file{Single Step} \\
    \addlinespace[2pt]
    \eng{Breakpoint} & \code{G} pada alamat; jumlah terbatas & Titik
    henti pada baris sumber & Titik henti pada baris sumber \\
    \addlinespace[2pt]
    Memeriksa memori & \code{D} dan \code{E} & Jendela \eng{dump}
    memori & Jendela memori \\
    \addlinespace[2pt]
    Memeriksa \eng{stack} & \code{D SS:SP} & Jendela \eng{stack} &
    Jendela \eng{stack} \\
    \addlinespace[2pt]
    Memodifikasi kode mesin & \code{A} dan \code{E} & Tidak langsung
    & Tidak langsung \\
    \addlinespace[2pt]
    Menjalankan \file{.exe} lama & Ya & Ya & Perlu emulasi \\
    \bottomrule
\end{tabular}
\end{center}
\Sumber{Bab XV, Tabel perbandingan alat debugging}
\end{frame}

% ---------------------------------------------------------
\begin{frame}{Rangkuman Bagian 1}
\small
\begin{itemize}
    \item \eng{Debugging} adalah proses sistematis yang membuktikan
          penyebab kesalahan melalui bukti pengamatan pada level
          instruksi --- bukan kegiatan menebak-nebak.

    \vspace{0.25em}
    \item Lima jenis kesalahan adalah sintaks, saat berjalan, logika,
          segmentasi, dan \eng{overflow} (aritmatika maupun
          \eng{stack}).

    \vspace{0.25em}
    \item Hanya kesalahan sintaks yang dapat ditemukan tanpa
          \eng{debugger}; kesalahan logika paling sulit karena
          program selesai normal tetapi hasilnya salah.

    \vspace{0.25em}
    \item Tiga alat \eng{debugging} --- MS-DEBUG, GUI Turbo
          Assembler, dan Emu8086 --- menampilkan informasi yang sama;
          MS-DEBUG dipakai sebagai acuan karena memperlihatkan
          keadaan mesin apa adanya.
\end{itemize}
\end{frame}

% =========================================================
% BAGIAN 2
% =========================================================
\PemisahBagian{Menelusuri Program dengan Perintah DEBUG}

% ---------------------------------------------------------
\begin{frame}[fragile]{Memuat dan Menulis Berkas: \code{N}, \code{L}, \code{W}}
\small
\code{N} (\eng{name}) menetapkan nama berkas, \code{L} (\eng{load})
memuat isi berkas ke memori, dan \code{W} (\eng{write}) menulis
kembali isi memori ke berkas.

\vspace{0.2em}
\begin{keluaran}
C:\ASM>debug bug-loop.exe
-n bug-loop.exe
-l
\end{keluaran}

\vspace{0.2em}
\small
\begin{itemize}
    \item \code{L} \emph{tidak} menjalankan program; ia hanya
          memindahkan isi berkas ke memori dan menetapkan register
          segmen beserta titik masuknya.
    \item \code{L} dan \code{W} bekerja pada berkas yang
          \emph{terakhir} ditetapkan \code{N}, sehingga urutan
          \code{N} lalu \code{L} harus selalu dijaga.
\end{itemize}
\Sumber{Bab XV, Memuat dan Menulis Berkas}
\end{frame}

% ---------------------------------------------------------
\begin{frame}[fragile]{Perintah \code{R}: Memeriksa Register dan \eng{Flag}}
\small
\code{R} menampilkan seluruh register umum, register segmen, dan
seluruh \eng{flag} dalam satu layar. Baris kedua keluarannya juga
menampilkan instruksi yang sedang ditunjuk \reg{CS}:\reg{IP}.
\code{R} dapat pula dipakai untuk \emph{mengubah} isi register,
misalnya \code{R AX} atau \code{R IP}.

\vspace{0.2em}
\begin{keluaran}
-r
AX=13A6  BX=0000  CX=00FF  DX=0000  SP=0100  BP=0000  SI=0000  DI=0000
DS=1391  ES=1391  SS=13A7  CS=13A1  IP=0000   NV UP EI PL NZ NA PO NC
13A1:0000 B8A613         MOV     AX,13A6
\end{keluaran}

\vspace{0.2em}
{\footnotesize Delapan penanda \eng{flag} ditampilkan dengan urutan
yang selalu tetap: \flag{OF}, \flag{DF}, \flag{IF}, \flag{SF},
\flag{ZF}, \flag{AF}, \flag{PF}, lalu \flag{CF}. Menghafal urutan ini
sangat membantu ketika mengisi tabel \eng{trace}.\par}
\Sumber{Bab XV, Memeriksa Register dengan R}
\end{frame}

% ---------------------------------------------------------
\begin{frame}[fragile]{Perintah \code{U}: Membongkar Kode Mesin}
\small
\code{U} (\eng{unassemble}) menerjemahkan kembali kode mesin menjadi
mnemonic assembly. Perintah ini dipakai untuk memastikan bahwa kode
yang dijalankan benar-benar sama dengan yang ditulis, dan untuk
mengetahui di alamat berapa setiap instruksi berada.

\vspace{0.15em}
\begin{keluaran}[basicstyle=\ttfamily\fontsize{7}{8.4}\selectfont]
-u 0000 001c
13A1:0000 B8A613         MOV     AX,13A6
13A1:0003 8ED8           MOV     DS,AX
13A1:0005 31C0           XOR     AX,AX
13A1:0007 B90500         MOV     CX,0005
13A1:000A 01C8           ADD     AX,CX
13A1:000C 41             INC     CX
13A1:000D E2FB           LOOP    000A
13A1:000F B409           MOV     AH,09
13A1:0011 BA0000         MOV     DX,0000
13A1:0014 CD21           INT     21
13A1:0016 B44C           MOV     AH,4C
13A1:0018 B000           MOV     AL,00
13A1:001A CD21           INT     21
\end{keluaran}
\Sumber{Bab XV, Membongkar Kode Mesin dengan U}
\end{frame}

% ---------------------------------------------------------
\begin{frame}[fragile]{Penelusuran Satu Langkah dengan \code{T}}
\small
\code{T} (\eng{trace}) menjalankan \emph{tepat satu} instruksi, lalu
menampilkan kembali seluruh register dan \eng{flag}, serta mnemonic
instruksi berikutnya. Inilah perintah terpenting pada bab ini: hanya
dengan menjalankan satu instruksi demi satu instruksi kita dapat
melihat pada langkah berapa sebuah nilai mulai menyimpang.

\vspace{0.2em}
\begin{keluaran}
-t
AX=13A6  BX=0000  CX=00FF  DX=0000  SP=0100  BP=0000  SI=0000  DI=0000
DS=13A6  ES=1391  SS=13A7  CS=13A1  IP=0005   NV UP EI PL NZ NA PO NC
13A1:0005 31C0           XOR     AX,AX
\end{keluaran}

\vspace{0.2em}
{\footnotesize\color{obeGray}Baris terakhir selalu menunjukkan
instruksi yang \emph{akan} dijalankan berikutnya, bukan yang baru
saja selesai. Pada langkah di atas \reg{DS} sudah bernilai \hex{13A6}
karena \asm{MOV DS, AX} baru saja dieksekusi.\par}
\Sumber{Bab XV, Penelusuran Satu Langkah dengan T}
\end{frame}

% ---------------------------------------------------------
\begin{frame}[fragile]{Sesi Perulangan: \reg{CX} yang Tidak Pernah Nol}
\scriptsize Inti temuan pada \file{bug-loop.exe}: \reg{CX} naik
menjadi \hex{0006} oleh \asm{INC CX}, lalu \asm{LOOP} menguranginya
kembali menjadi \hex{0005}.

\vspace{0.1em}
\begin{keluaran}[basicstyle=\ttfamily\fontsize{7}{8.4}\selectfont]
-t
AX=0000  BX=0000  CX=0005  DX=0000  SP=0100  BP=0000  SI=0000  DI=0000
DS=13A6  ES=1391  SS=13A7  CS=13A1  IP=000A   NV UP EI PL ZR NA PE NC
13A1:000A 01C8           ADD     AX,CX
-t
AX=0005  BX=0000  CX=0005  DX=0000  SP=0100  BP=0000  SI=0000  DI=0000
13A1:000C 41             INC     CX
-t
AX=0005  BX=0000  CX=0006  DX=0000  SP=0100  BP=0000  SI=0000  DI=0000
13A1:000D E2FB           LOOP    000A
-t
AX=0005  BX=0000  CX=0005  DX=0000  SP=0100  BP=0000  SI=0000  DI=0000
13A1:000A 01C8           ADD     AX,CX
\end{keluaran}
\Sumber{Bab XV, Sesi DEBUG lengkap pada bug-loop.exe}
\end{frame}

% ---------------------------------------------------------
\begin{frame}{Perintah DEBUG Inti: Penelusuran}
\scriptsize
\begin{center}
\begin{tabular}{@{}p{1.2cm}p{3.3cm}p{2.5cm}p{5.4cm}@{}}
    \toprule
    \textbf{Perintah} & \textbf{Kegunaan} & \textbf{Contoh} &
    \textbf{Keluaran yang muncul} \\
    \midrule
    \code{R} & Menampilkan dan mengubah register serta \eng{flag} &
    \code{R} atau \code{R AX} &
    \code{AX=13A6 ... NV UP EI PL NZ NA PO NC} \\
    \addlinespace[2pt]
    \code{U} & Membongkar kode mesin menjadi mnemonic &
    \code{U 0000 001C} &
    \code{13A1:0000 B8A613 MOV AX,13A6} \\
    \addlinespace[2pt]
    \code{T} & Menjalankan satu instruksi, termasuk masuk ke dalam
    \asm{CALL} & \code{T} &
    Register baru, lalu mnemonic berikutnya \\
    \addlinespace[2pt]
    \code{P} & Menjalankan sampai instruksi berikutnya, melewati
    \asm{CALL}/\code{INT} & \code{P} &
    Register baru setelah prosedur selesai \\
    \addlinespace[2pt]
    \code{G} & Menjalankan sampai alamat tertentu (\eng{breakpoint})
    & \code{G 000A} &
    Berhenti di \hex{000A} dan menampilkan register \\
    \addlinespace[2pt]
    \code{D} & Menampilkan isi memori (\eng{dump}) & \code{D DS:0} &
    \code{13A6:0000 4A 75 6D 6C ... Jumlah 1..5 = \$.} \\
    \bottomrule
\end{tabular}
\end{center}
\Sumber{Bab XV, Tabel perintah DEBUG}
\end{frame}

% ---------------------------------------------------------
\begin{frame}{Perintah DEBUG Inti: Memori dan Berkas}
\scriptsize
\begin{center}
\begin{tabular}{@{}p{1.2cm}p{3.3cm}p{2.5cm}p{5.4cm}@{}}
    \toprule
    \textbf{Perintah} & \textbf{Kegunaan} & \textbf{Contoh} &
    \textbf{Keluaran yang muncul} \\
    \midrule
    \code{E} & Mengubah isi memori satu per satu &
    \code{E DS:0 48 61 6C 6F 24} &
    Memori \reg{DS}:\hex{0000} berubah menjadi \code{Halo\$} \\
    \addlinespace[2pt]
    \code{A} & Merakit instruksi langsung di memori &
    \code{A 0100} &
    \code{13A1:0100 MOV AH,02} (DEBUG meminta baris berikutnya) \\
    \addlinespace[2pt]
    \code{N} & Menetapkan nama berkas & \code{N latihan.com} &
    Tidak ada keluaran \\
    \addlinespace[2pt]
    \code{L} & Memuat berkas yang telah dinamai ke memori &
    \code{L} & Tidak ada keluaran; register terisi \\
    \addlinespace[2pt]
    \code{W} & Menulis isi memori ke berkas & \code{W 0100} &
    \code{Writing 0020 bytes} \\
    \bottomrule
\end{tabular}
\end{center}

\vspace{0.2em}
{\scriptsize\color{obeGray}\code{L} dan \code{W} selalu bekerja pada
berkas yang terakhir ditetapkan \code{N}; urutan \code{N} lalu
\code{L} tidak boleh dibalik.\par}
\Sumber{Bab XV, Tabel perintah DEBUG}
\end{frame}

% ---------------------------------------------------------
\begin{frame}[fragile]{Perbedaan \code{T} dan \code{P}}
\small
Keduanya sama-sama menjalankan program, tetapi berbeda dalam
menyikapi instruksi yang memindahkan kendali.

\vspace{0.25em}
\begin{columns}[t]
\begin{column}{0.48\textwidth}
\begin{exampleblock}{\code{T} --- \eng{trace}}
\footnotesize Menjalankan satu instruksi \emph{apa pun}. Jika
instruksi itu \asm{CALL}, \code{T} \emph{masuk} ke dalam prosedur;
jika \code{INT 21h}, \code{T} masuk ke rutin layanan DOS.
\end{exampleblock}
\end{column}
\begin{column}{0.48\textwidth}
\begin{alertblock}{\code{P} --- \eng{proceed}}
\footnotesize Menjalankan instruksi sampai instruksi berikutnya pada
level yang sama. Jika instruksi itu \asm{CALL} atau \code{INT},
\code{P} menjalankannya sampai selesai lalu berhenti setelahnya.
\end{alertblock}
\end{column}
\end{columns}

\vspace{0.25em}
{\footnotesize\important{Pakai \code{P} untuk \code{INT 21h} dan
prosedur yang sudah diyakini benar; pakai \code{T} untuk prosedur
yang sedang dicurigai.} \code{P} \emph{tidak memperbaiki} bug di
dalam prosedur --- ia hanya menyembunyikan langkah-langkahnya.\par}
\Sumber{Bab XV, Perbedaan T dan P}
\end{frame}

% ---------------------------------------------------------
\begin{frame}{Breakpoint dengan Perintah \code{G}}
\small
Menelusuri ratusan instruksi dengan \code{T} jelas tidak efisien.
Perintah \code{G} (\eng{go}) menjalankan program sampai alamat
tertentu berhenti, sehingga programmer dapat langsung melompat ke
bagian yang dicurigai.

\vspace{0.2em}
\begin{itemize}
    \item \code{G 000A} menjalankan instruksi dari \hex{0000} sampai
          \hex{000A} dan berhenti \emph{sebelum} instruksi pada
          alamat itu dieksekusi.
    \item Beberapa alamat dapat ditulis sekaligus, misalnya
          \code{G 000A 000F}; program berhenti pada alamat pertama
          yang tercapai, sehingga titik di dalam dan sesudah
          perulangan dapat diperiksa dalam satu kali jalan.
\end{itemize}

\vspace{0.25em}
\begin{alertblock}{Keterbatasan DEBUG klasik}
\footnotesize (1) jumlah \eng{breakpoint} terbatas, umumnya empat
sampai sepuluh; (2) \eng{breakpoint} hanya pada \emph{alamat}, bukan
baris sumber atau kondisi seperti ``berhenti bila \reg{AX} bernilai
nol''; (3) hanya mode real 16-bit, sehingga program 32-bit tidak
dapat ditelusuri.
\end{alertblock}
\Sumber{Bab XV, Breakpoint dengan Perintah G}
\end{frame}

% ---------------------------------------------------------
\begin{frame}{Tabel Trace Register dan \eng{Flag}}
\scriptsize
\begin{center}
\begin{tabular}{@{}p{0.6cm}p{1.6cm}p{2.3cm}p{1.0cm}p{1.0cm}p{1.0cm}p{0.6cm}p{0.6cm}p{0.6cm}@{}}
    \toprule
    \textbf{L} & \textbf{Alamat} & \textbf{Instruksi} & \textbf{AX} &
    \textbf{CX} & \textbf{SP} & \textbf{ZF} & \textbf{CF} &
    \textbf{OF} \\
    \midrule
    1 & 13A1:0000 & MOV AX,13A6 & 13A6 & 0000 & 0100 & 0 & 0 & 0 \\
    \addlinespace[2pt]
    2 & 13A1:0003 & MOV DS,AX & 13A6 & 0000 & 0100 & 0 & 0 & 0 \\
    \addlinespace[2pt]
    3 & 13A1:0005 & XOR AX,AX & 0000 & 0000 & 0100 & 1 & 0 & 0 \\
    \addlinespace[2pt]
    4 & 13A1:0007 & MOV CX,0005 & 0000 & 0005 & 0100 & 1 & 0 & 0 \\
    \addlinespace[2pt]
    5 & 13A1:000A & ADD AX,CX & 0005 & 0005 & 0100 & 0 & 0 & 0 \\
    \addlinespace[2pt]
    6 & 13A1:000C & INC CX & 0005 & 0006 & 0100 & 0 & 0 & 0 \\
    \addlinespace[2pt]
    7 & 13A1:000D & LOOP 000A & 0005 & 0005 & 0100 & 0 & 0 & 0 \\
    \addlinespace[2pt]
    8 & 13A1:000A & ADD AX,CX & 000A & 0005 & 0100 & 0 & 0 & 0 \\
    \addlinespace[2pt]
    9 & 13A1:000C & INC CX & 000A & 0006 & 0100 & 0 & 0 & 0 \\
    \addlinespace[2pt]
    10 & 13A1:000D & LOOP 000A & 000A & 0005 & 0100 & 0 & 0 & 0 \\
    \addlinespace[2pt]
    11 & 13A1:000A & ADD AX,CX & 000F & 0005 & 0100 & 0 & 0 & 0 \\
    \bottomrule
\end{tabular}
\end{center}

\vspace{0.1em}
{\scriptsize\color{obeGray}Penyebabnya tampak pada langkah 6 dan 7:
\reg{CX} berayun pada \dec{5} dan \dec{6} sehingga tidak pernah
mencapai nol. \flag{ZF} pada langkah 4 masih \dec{1} karena \asm{MOV}
\emph{tidak pernah} mengubah \eng{flag}.\par}
\Sumber{Bab XV, Tabel trace register dan flag}
\end{frame}

% ---------------------------------------------------------
\begin{frame}{Rangkuman Bagian 2}
\small
\begin{itemize}
    \item Berkas dimuat dengan urutan \code{N} lalu \code{L};
          \code{L} hanya memindahkan isi berkas ke memori, bukan
          menjalankannya.

    \vspace{0.22em}
    \item \code{R} menampilkan register dan \eng{flag} dengan urutan
          tetap \flag{OF}--\flag{DF}--\flag{IF}--\flag{SF}--\flag{ZF}
          --\flag{AF}--\flag{PF}--\flag{CF}, sekaligus instruksi
          berikutnya; \code{U} membongkar kode mesin menjadi mnemonic.

    \vspace{0.22em}
    \item \code{T} menjalankan satu instruksi dan masuk ke dalam
          \asm{CALL} serta \code{INT}; \code{P} melewatinya.
          \code{P} tidak memperbaiki bug, hanya menyembunyikan
          langkahnya.

    \vspace{0.22em}
    \item \code{G alamat} memasang \eng{breakpoint}, tetapi DEBUG
          klasik hanya menyediakan titik henti dalam jumlah terbatas
          dan hanya pada alamat, bukan pada baris sumber atau
          kondisi.

    \vspace{0.22em}
    \item Bukti utama \eng{debugging} adalah tabel \eng{trace}: pada
          setiap langkah dicatat alamat, instruksi, register, dan
          \eng{flag}.
\end{itemize}
\end{frame}

% =========================================================
% BAGIAN 3
% =========================================================
\PemisahBagian{Analisis Memori dan Stack serta Metodologi Debugging}

% ---------------------------------------------------------
\begin{frame}[fragile]{Membaca Keluaran \code{D}}
\small
\code{D} menampilkan \dec{128} byte memori dalam delapan baris.
Setiap baris memuat tiga bagian yang dibaca secara berbeda.

\vspace{0.2em}
\begin{keluaran}
-d ds:0
13A6:0000 4A 75 6D 6C 61 68 20 31-2E 2E 35 20 3D 20 24 00  Jumlah 1..5 = $.
\end{keluaran}

\vspace{0.2em}
\small
\begin{itemize}
    \item \textbf{Alamat} dalam bentuk \eng{segmen}:\eng{offset},
          yaitu \hex{13A6}:\hex{0000}.
    \item \textbf{Isi memori}: enam belas pasang digit heksadesimal;
          setiap dua digit adalah satu byte.
    \item \textbf{Kolom karakter ASCII} di sebelah kanan; byte yang
          bukan karakter cetak ditampilkan sebagai titik.
\end{itemize}

\vspace{0.15em}
{\footnotesize\color{obeGray}Kolom karakter memperlihatkan string
\code{Jumlah 1..5 = \$}; kolom heksadesimal memperlihatkan kode
ASCII-nya, \hex{4A} untuk \code{J}, \hex{75} untuk \code{u}, sampai
\hex{24} untuk tanda \code{\$}; byte terakhir \hex{00} adalah penanda
akhir yang dipasang assembler.\par}
\Sumber{Bab XV, Membaca Keluaran D}
\end{frame}

% ---------------------------------------------------------
\begin{frame}[fragile]{Memeriksa Variabel: Kesalahan \eng{Off-by-One}}
\small
Program bermaksud menjumlahkan lima byte larik \code{nilai}; jumlah
yang benar adalah \dec{150} (\hex{0096}).

\vspace{0.15em}
\begin{keluaran}
-d ds:0
13A6:0000  0A 14 1E 28 32 B7 00 00-00 00 00 00 00 00 00 00  ...(2...........
\end{keluaran}

\vspace{0.15em}
{\footnotesize
\begin{itemize}
    \item Lima byte pertama sesuai data: \hex{0A}=\dec{10} sampai
          \hex{32}=\dec{50}.
    \item Byte keenam bernilai \hex{B7}, yaitu sisa data lama di
          memori, bukan bagian dari larik.
    \item Karena penghitung bernilai \dec{6}, byte itu ikut
          dijumlahkan sehingga akumulator menghasilkan \hex{014D}
          (\dec{333}) alih-alih \hex{0096}.
\end{itemize}}

\vspace{0.1em}
{\footnotesize\color{obeGray}Bukti ini mengubah dugaan menjadi
kepastian: kesalahannya bukan pada penjumlahan, melainkan pada batas
perulangan.\par}
\Sumber{Bab XV, Memeriksa Variabel dengan D}
\end{frame}

% ---------------------------------------------------------
\begin{frame}[fragile]{Program Larik Ber-kesalahan \eng{Off-by-One}}
\scriptsize
\begin{columns}[t]
\begin{column}{0.62\textwidth}
\begin{kodeasm}[numbers=none, basicstyle=\ttfamily\fontsize{7.5}{8.4}\selectfont]
.DATA
    nilai  DB  10, 20, 30, 40, 50   ; lima byte
.CODE
START:
    XOR  DX, DX          ; DX = akumulator jumlah
    MOV  CX, 6           ; BUG: seharusnya MOV CX, 5
    MOV  SI, 0           ; SI = indeks larik
ULANG:
    MOV  AH, 0           ; nolkan AH
    MOV  AL, nilai[SI]   ; ambil elemen ke-SI
    ADD  DX, AX          ; tambahkan ke akumulator
    INC  SI              ; naikkan indeks
    LOOP ULANG           ; ulangi sebanyak CX
\end{kodeasm}
\end{column}
\begin{column}{0.35\textwidth}
{\footnotesize
\begin{itemize}
    \item Larik berisi lima byte, sedangkan penghitung diisi
          \dec{6}, sehingga satu byte di luar larik ikut dibaca.
    \item Perbaikannya: ubah \asm{MOV CX, 6} menjadi
          \asm{MOV CX, 5}.
    \item Cara yang lebih aman adalah menghitung jumlah elemen
          otomatis dengan \code{panjang EQU \$ - nilai}.
\end{itemize}}
\end{column}
\end{columns}

\vspace{0.15em}
{\scriptsize\color{obeGray}Program ditutup \asm{MOV AH, 4Ch} dan
\code{INT 21h}. Sumber: Bab XV, Studi Kasus Tiga Bug (Kasus 2).\par}
\end{frame}

% ---------------------------------------------------------
\begin{frame}{Membaca Bingkai Stack}
\begin{center}
\begin{tikzpicture}[
    sel/.style={draw=obeBlue,line width=0.7pt,fill=obeBlueLight,
        minimum width=5.6cm,minimum height=0.62cm,font=\small,
        inner sep=2pt},
    selsp/.style={draw=obeMaroon,line width=1.0pt,
        fill=obeMaroonLight,minimum width=5.6cm,minimum height=0.62cm,
        font=\small,inner sep=2pt},
    al/.style={font=\scriptsize\ttfamily,text=obeGray,anchor=east},
    sp/.style={-{Stealth[length=2.5mm]},line width=1.0pt,obeMaroon}
]
    \node[sel,fill=obeGrayLight] (d) at (0,0) {dasar \eng{stack}};
    \node[sel,below=0.08cm of d] (r) {alamat kembali \hex{0008}};
    \node[sel,below=0.08cm of r] (ax) {AX lama = \hex{13A6}};
    \node[sel,below=0.08cm of ax] (bx) {BX lama = \hex{0000}};
    \node[selsp,below=0.08cm of bx] (cx) {CX lama = \hex{00FF}};

    \node[al] at ([xshift=-0.18cm]d.west) {\hex{0100}};
    \node[al] at ([xshift=-0.18cm]r.west) {\hex{00FE}};
    \node[al] at ([xshift=-0.18cm]ax.west) {\hex{00FC}};
    \node[al] at ([xshift=-0.18cm]bx.west) {\hex{00FA}};
    \node[al] at ([xshift=-0.18cm]cx.west) {\hex{00F8}};

    \draw[sp] ([xshift=0.18cm]cx.east) -- ++(0.55cm,0)
        node[right,font=\scriptsize,text=obeMaroon,align=left]
        {\reg{SP} menunjuk\\puncak \eng{stack}};
\end{tikzpicture}
\end{center}

\vspace{0.15em}
{\footnotesize\color{obeGray}\asm{CALL} menurunkan \reg{SP} dari
\hex{0100} ke \hex{00FE} dan menuliskan alamat kembali di sana;
ketiga \asm{PUSH} menurunkannya sampai \hex{00F8}. Alamat kembali
berada \emph{di bawah} seluruh data yang ditekan prosedur, sehingga
\asm{RET} hanya dapat mengambilnya bila seluruh data di atasnya
sudah diambil kembali. Sumber: Bab XV, Membaca Bingkai Stack.\par}
\end{frame}

% ---------------------------------------------------------
\begin{frame}[fragile]{Menemukan Stack Tidak Seimbang}
\scriptsize Prosedur menekan tiga register tetapi hanya memulihkan
dua, sehingga \asm{RET} membaca nilai yang bukan alamat kembali.

\vspace{0.1em}
\begin{keluaran}[basicstyle=\ttfamily\fontsize{7}{8.4}\selectfont]
-g 0005
13A1:0005 E80D00         CALL    0015
-t
AX=13A6  BX=0000  CX=00FF  DX=0000  SP=00FE  BP=0000  SI=0000  DI=0000
13A1:0015 50             PUSH    AX
-t
AX=13A6  BX=0000  CX=00FF  DX=0000  SP=00FC  BP=0000  SI=0000  DI=0000
13A1:0016 53             PUSH    BX
-t
AX=13A6  BX=0000  CX=00FF  DX=0000  SP=00FA  BP=0000  SI=0000  DI=0000
13A1:0017 51             PUSH    CX
-d ss:00f8
13A7:00F8  FF 00 00 00 A6 13 08 00-00 00 00 00 00 00 00 00  ................
-t
AX=0003  BX=0002  CX=00FF  DX=0000  SP=00FC  BP=0000  SI=0000  DI=0000
13A1:0024 C3             RET
\end{keluaran}

\vspace{0.1em}
{\scriptsize\color{obeGray}Setelah dua \asm{POP}, \reg{SP} berhenti
pada \hex{00FC} dan menunjuk \code{A6 13}, yaitu nilai \reg{AX} lama;
\asm{RET} memuatnya ke \reg{IP} sehingga program melompat ke
\hex{13A6} --- jauh di luar batas program.\par}
\Sumber{Bab XV, Menemukan stack tidak seimbang melalui D dan T}
\end{frame}

% ---------------------------------------------------------
\begin{frame}{Deteksi \eng{Overflow} melalui \eng{Flag} dan \reg{SP}}
\small
\eng{Overflow} aritmatika tidak selalu menghentikan program, tetapi
selalu meninggalkan jejak pada \eng{flag}.

\vspace{0.3em}
\begin{columns}[t]
\begin{column}{0.48\textwidth}
\begin{block}{\eng{Overflow} aritmatika}
\footnotesize Operasi yang melewati batas bilangan
\emph{bertanda} menyalakan \flag{OF}, sedangkan yang melewati batas
bilangan \emph{tak bertanda} menyalakan \flag{CF}.
\end{block}
\end{column}
\begin{column}{0.48\textwidth}
\begin{alertblock}{\eng{Overflow} \eng{stack}}
\footnotesize Meninggalkan jejak berbeda: nilai \reg{SP} terus
menurun tanpa pernah naik kembali. Gejalanya adalah perubahan nilai
variabel yang tidak pernah disentuh program.
\end{alertblock}
\end{column}
\end{columns}

\vspace{0.3em}
{\footnotesize\color{obeGray}Karena itu \flag{CF}, \flag{OF}, dan
register \reg{SP} sebaiknya selalu disertakan pada tabel
\eng{trace}.\par}
\Sumber{Bab XV, Deteksi Overflow melalui Flag dan SP}
\end{frame}

% ---------------------------------------------------------
\begin{frame}{Metodologi Debugging yang Sistematis}
\begin{center}
\begin{tikzpicture}[
    tahap/.style={draw=obeTeal,line width=0.8pt,fill=obeTealLight,
        rounded corners=3pt,text width=3.4cm,minimum height=1.25cm,
        align=center,inner sep=3pt,font=\footnotesize},
    panah/.style={-{Stealth[length=2.2mm]},line width=0.9pt,obeTeal}
]
    \node[tahap] (t1) {\textbf{1. Reproduksi}\\{\scriptsize\itshape
        pastikan dapat diulang}};
    \node[tahap,right=0.5cm of t1] (t2) {\textbf{2. Isolasi}\\
        {\scriptsize\itshape perkecil cakupan}};
    \node[tahap,right=0.5cm of t2] (t3) {\textbf{3. Hipotesis}\\
        {\scriptsize\itshape dugaan yang dapat diuji}};
    \node[tahap,below=0.85cm of t1] (t4) {\textbf{4. Uji hipotesis}\\
        {\scriptsize\itshape buktikan dengan \code{T}, \code{G},
        \code{D}}};
    \node[tahap,right=0.5cm of t4] (t5) {\textbf{5. Perbaiki}\\
        {\scriptsize\itshape ubah satu hal saja}};
    \node[tahap,right=0.5cm of t5] (t6) {\textbf{6. Verifikasi}\\
        {\scriptsize\itshape ulangi seluruh pengujian}};

    \draw[panah] (t1) -- (t2);
    \draw[panah] (t2) -- (t3);
    \draw[panah] (t3.south) -- ++(0,-0.42) -| (t4.north);
    \draw[panah] (t4) -- (t5);
    \draw[panah] (t5) -- (t6);
\end{tikzpicture}
\end{center}

\vspace{0.15em}
{\footnotesize\color{obeGray}Pada tahap perbaikan, ubah \emph{satu
hal saja} agar penyebab keberhasilan dapat dipastikan. Bila program
kemudian berjalan benar, mengubah beberapa baris sekaligus membuat
kita tidak pernah tahu baris mana yang sebenarnya
memperbaiki keadaan.\par}
\Sumber{Bab XV, Metodologi Debugging yang Sistematis}
\end{frame}

% ---------------------------------------------------------
\begin{frame}{Studi Kasus Tiga Bug}
\footnotesize
\begin{enumerate}
    \item \term{Perulangan tak berujung.} \file{bug-loop.asm} berhenti
          tidak pada waktunya karena \reg{CX} dinaikkan di dalam
          badan perulangan; \asm{LOOP} menguranginya kembali,
          sehingga \reg{CX} berayun pada \dec{5} dan \dec{6}.
          Perbaikannya menghapus \asm{INC CX}.

    \vspace{0.25em}
    \item \term{Kesalahan \eng{off-by-one} pada larik.} Program
          selesai normal, sehingga kesalahan sintaks dan kesalahan
          saat berjalan dapat disingkirkan. \code{D DS:0}
          memperlihatkan byte keenam di luar larik; perbaikannya
          mengubah \asm{MOV CX, 6} menjadi \asm{MOV CX, 5}.

    \vspace{0.25em}
    \item \term{\eng{Stack} tidak seimbang.} \file{bug-stack.asm}
          melompat ke alamat aneh --- gejala khas \eng{stack} rusak,
          karena hanya \asm{CALL}, \asm{RET}, dan \code{INT} yang
          dapat memindahkan kendali ke luar batas program. Prosedur
          menekan tiga register tetapi memulihkan dua; perbaikannya
          menambahkan satu \asm{POP AX}.
\end{enumerate}

\vspace{0.2em}
{\footnotesize\color{obeGray}Pola yang sama pada ketiganya:
gejalanya muncul di satu tempat, sedangkan penyebabnya berada di
tempat lain.\par}
\Sumber{Bab XV, Studi Kasus Tiga Bug}
\end{frame}

% ---------------------------------------------------------
\begin{frame}{Rangkuman Bagian 3}
\small
\begin{itemize}
    \item Keluaran \code{D} dibaca dalam tiga bagian: alamat, isi
          memori heksadesimal, dan kolom karakter ASCII. Kolom
          karakter paling cepat dipakai untuk mengenali letak
          sebuah string.

    \vspace{0.22em}
    \item Bagian \eng{stack} dibaca dari alamat terendah ke
          tertinggi: data yang ditekan terakhir berada pada alamat
          terendah, sedangkan alamat kembali berada di bawah seluruh
          data tersebut.

    \vspace{0.22em}
    \item \eng{Overflow} aritmatika meninggalkan jejak pada
          \flag{OF} (bilangan bertanda) dan \flag{CF} (bilangan tak
          bertanda); \eng{overflow} \eng{stack} meninggalkan
          \reg{SP} yang terus menurun.

    \vspace{0.22em}
    \item Metodologi \eng{debugging} terdiri atas enam tahap:
          reproduksi, isolasi, hipotesis, pengujian, perbaikan, dan
          verifikasi --- dengan perbaikan satu hal pada satu waktu.

    \vspace{0.22em}
    \item Tiga bug yang dibedah adalah perulangan tak berujung,
          kesalahan \eng{off-by-one} pada larik, dan \eng{stack}
          tidak seimbang.
\end{itemize}
\end{frame}

% =========================================================
% CONTOH PROGRAM
% =========================================================
\begin{frame}[fragile]{Contoh 1: Perulangan Tak Berujung}
\scriptsize Bermaksud menjumlahkan bilangan 1 sampai 5 (seharusnya
\hex{000F} atau \dec{15}); program dirakit tanpa kesalahan sintaks,
tetapi tidak pernah selesai.

\vspace{0.1em}
\begin{columns}[t]
\begin{column}{0.5\textwidth}
{\footnotesize\color{obeMaroon}\bfseries \file{bug-loop.asm}
(ber-bug)}\\[0.1em]
\begin{kodeasm}[numbers=none, basicstyle=\ttfamily\fontsize{7.5}{8.4}\selectfont]
XOR AX, AX      ; akumulator = 0
MOV CX, 5       ; penghitung = 5
ULANG:
    ADD AX, CX  ; jumlahkan penghitung
    INC CX      ; BUG: CX dinaikkan
    LOOP ULANG  ; LOOP mengurangi CX
MOV AH, 09h     ; tak pernah dicapai
LEA DX, pesan
INT 21h
\end{kodeasm}
\end{column}
\begin{column}{0.46\textwidth}
{\footnotesize\color{obeTeal}\bfseries \file{bug-loop-fix.asm}
(perbaikan)}\\[0.1em]
\begin{kodeasm}[numbers=none, basicstyle=\ttfamily\fontsize{7.5}{8.4}\selectfont]
XOR AX, AX      ; akumulator = 0
MOV CX, 5       ; penghitung = 5
ULANG:
    ADD AX, CX
    LOOP ULANG  ; CX turun otomatis
MOV BX, AX      ; BX = 15
MOV AH, 09h     ; cetak judul
LEA DX, pesan
INT 21h
\end{kodeasm}

\vspace{0.1em}
{\footnotesize\color{obeTeal}\bfseries Keluaran}\\[0.1em]
\begin{keluaran}
Jumlah 1..5 = 15
\end{keluaran}
\end{column}
\end{columns}
\end{frame}

% ---------------------------------------------------------
\begin{frame}[fragile]{Contoh 2: Stack yang Tidak Seimbang}
\scriptsize Program memanggil prosedur \asm{HITUNG} lalu bermaksud
mencetak pesan penutup; kendali melompat ke alamat yang tidak dikenal
dan pesan itu tidak pernah muncul.

\vspace{0.1em}
\begin{columns}[t]
\begin{column}{0.5\textwidth}
{\footnotesize\color{obeMaroon}\bfseries \file{bug-stack.asm}
(ber-bug)}\\[0.1em]
\begin{kodeasm}[numbers=none, basicstyle=\ttfamily\fontsize{7.5}{8.4}\selectfont]
HITUNG PROC
    PUSH AX     ; simpan AX lama
    PUSH BX     ; simpan BX lama
    PUSH CX     ; simpan CX lama
    MOV  AX, 0001h
    MOV  BX, 0002h
    ADD  AX, BX ; AX = 0003h
    MOV  CX, AX
    POP  CX     ; pulihkan CX
    POP  BX     ; pulihkan BX
    RET         ; BUG: POP AX tidak ada
HITUNG ENDP
\end{kodeasm}
\end{column}
\begin{column}{0.46\textwidth}
{\footnotesize\color{obeTeal}\bfseries \file{bug-stack-fix.asm}
(perbaikan)}\\[0.1em]
\begin{kodeasm}[numbers=none, basicstyle=\ttfamily\fontsize{7.5}{8.4}\selectfont]
    MOV  CX, AX
    POP  CX     ; pulihkan CX
    POP  BX     ; pulihkan BX
    POP  AX     ; PERBAIKAN: pulihkan AX
    RET         ; alamat kembali benar
HITUNG ENDP
\end{kodeasm}

\vspace{0.1em}
{\footnotesize\color{obeTeal}\bfseries Keluaran}\\[0.1em]
\begin{keluaran}
Selesai.
\end{keluaran}

\vspace{0.1em}
{\footnotesize Setiap \asm{PUSH} memerlukan \asm{POP} pasangannya
pada jalur eksekusi yang sama.\par}
\end{column}
\end{columns}
\end{frame}

% =========================================================
% PRAKTIKUM
% =========================================================
\begin{frame}{Praktikum: Debugging Terpandu Tiga Kasus Bug}
\begin{block}{Judul praktikum}
\small \eng{Debugging} terpandu tiga kasus bug tingkat instruksi
\quad\textbullet\quad \textbf{Sub-CPMK:} 4.1 dan 4.2
\quad\textbullet\quad \textbf{Waktu:} 100 menit
\quad\textbullet\quad \textbf{Bentuk:} individu, laporan tertulis
\end{block}

\vspace{0.1em}
\begin{columns}[t]
\begin{column}{0.44\textwidth}
{\color{obeBlue}\bfseries Alat dan bahan}
\begin{itemize}\footnotesize
    \item Komputer bersistem operasi Windows
    \item DOSBox beserta berkas \file{DEBUG.EXE}
    \item TASM dan TLINK (atau Emu8086 sebagai pembanding)
    \item Berkas \file{bug-loop.asm}, \file{bug-stack.asm}, dan
          \file{program-larik.asm}
\end{itemize}
\end{column}
\begin{column}{0.52\textwidth}
{\color{obeBlue}\bfseries Langkah kerja}
\begin{enumerate}\footnotesize
    \item Rakit ketiga berkas uji dengan TASM dan TLINK.
    \item \textbf{Kasus 1:} muat \file{bug-loop.exe}, catat \code{R}
          dan \code{U}, lalu jalankan \code{T} sepuluh kali.
    \item \textbf{Kasus 2:} jalankan \code{G} sampai setelah
          \asm{LOOP}, lalu \code{D DS:0} untuk memeriksa larik.
    \item \textbf{Kasus 3:} \code{G 0005}, telusuri dengan \code{T},
          lalu bandingkan \code{T} dan \code{P} pada \asm{CALL}.
    \item Tuliskan hipotesis, perbaiki, rakit ulang, dan verifikasi.
\end{enumerate}
\end{column}
\end{columns}

\vspace{0.05em}
{\scriptsize\textbf{Dikumpulkan:} tabel \eng{trace} lengkap
(minimal sepuluh langkah per kasus) dengan kolom \reg{AX}, \reg{BX},
\reg{CX}, \reg{DX}, \reg{SP}, \flag{ZF}, \flag{CF}, \flag{OF};
salinan keluaran \code{R}, \code{U}, \code{D}, dan \code{T}; berkas
hasil perbaikan beserta tangkapan layarnya.\par}
\end{frame}

% =========================================================
% LATIHAN
% =========================================================
\begin{frame}{Latihan}
\small
\begin{enumerate}
    \item Jelaskan perbedaan lima jenis kesalahan program. Untuk
          setiap jenis, sebutkan satu contoh dan alat atau perintah
          yang paling tepat untuk menemukannya.

    \vspace{0.25em}
    \item Sebuah program memanggil \code{INT 21h} dengan fungsi
          \hex{09}. Apa yang terjadi bila ditelusuri dengan \code{T},
          dan apa yang terjadi bila ditelusuri dengan \code{P}?
          Mengapa \code{P} lebih dianjurkan pada kasus ini?

    \vspace{0.25em}
    \item Diberikan keluaran \code{R}: \code{AX=0000 CX=0005 SP=00F8
          IP=000A} dan baris \eng{flag} \code{NV UP EI NG NZ AC PO
          CY}. Nilai \flag{ZF}, \flag{CF}, \flag{SF}, dan \flag{OF}
          berturut-turut berapa? Instruksi apa yang dijalankan
          berikutnya?

    \vspace{0.25em}
    \item Diberikan keluaran \code{D DS:0000}: \code{48 41 4C 4F 24
          0A 0B ...}. String apa yang tersimpan di sana, dan pada
          byte berapa penanda akhirnya? Apa yang tampak di layar bila
          program membaca delapan byte mulai \hex{0000}?

    \vspace{0.25em}
    \item Tulislah prosedur \asm{HITUNG} yang menekan tiga register
          ke \eng{stack}, melakukan perhitungan sederhana, lalu
          memulihkan ketiganya sebelum \asm{RET}. Jelaskan mengapa
          urutan \asm{POP} harus kebalikan urutan \asm{PUSH}.
\end{enumerate}
\end{frame}

% =========================================================
% RANGKUMAN
% =========================================================
\begin{frame}{Rangkuman}
\footnotesize
\begin{itemize}
    \item \eng{Debugging} adalah proses sistematis untuk menemukan
          penyebab kesalahan melalui bukti pengamatan pada level
          instruksi, bukan melalui perubahan kode secara acak.

    \vspace{0.18em}
    \item Lima jenis kesalahan program: sintaks, saat berjalan,
          logika, segmentasi, dan \eng{overflow} (aritmatika maupun
          \eng{stack}). Hanya kesalahan sintaks yang dapat dideteksi
          tanpa \eng{debugger}.

    \vspace{0.18em}
    \item Tiga alat \eng{debugging} --- MS-DEBUG di DOSBox,
          \eng{debugger} terpadu GUI Turbo Assembler, dan jendela
          \eng{debug} Emu8086 --- menampilkan informasi yang sama:
          register, memori, \eng{flag}, dan \eng{stack}.

    \vspace{0.18em}
    \item Perintah wajib DEBUG adalah \code{R}, \code{U}, \code{T},
          \code{P}, \code{G}, \code{D}, \code{E}, \code{A}, serta
          \code{N}, \code{L}, dan \code{W}. Perintah \code{T} masuk
          ke dalam prosedur dan \code{INT}, sedangkan \code{P}
          melewatinya.

    \vspace{0.18em}
    \item Bukti utama \eng{debugging} adalah tabel \eng{trace}.
          Keluaran \code{D} dibaca dalam tiga bagian: alamat, isi
          heksadesimal, dan kolom karakter ASCII.

    \vspace{0.18em}
    \item Metodologi \eng{debugging} terdiri atas enam tahap:
          reproduksi, isolasi, hipotesis, pengujian, perbaikan, dan
          verifikasi --- dengan perbaikan satu hal pada satu waktu,
          lalu pengujian diulang pada seluruh kasus.
\end{itemize}
\end{frame}

% ---------------------------------------------------------
\begin{frame}{Cek Pemahaman}
\footnotesize
\begin{enumerate}
    \item Perintah DEBUG yang menampilkan seluruh register dan
          \eng{flag} adalah \dots
          \begin{enumerate}\footnotesize
              \item \code{U} \quad \item \code{D} \quad
              \item \code{R} \quad \item \code{T}
          \end{enumerate}

    \vspace{0.2em}
    \item Pada \asm{CALL HITUNG}, perintah yang \emph{tidak} masuk ke
          dalam prosedur adalah \dots
          \begin{enumerate}\footnotesize
              \item \code{T} \quad \item \code{P} \quad
              \item \code{U} \quad \item \code{R}
          \end{enumerate}

    \vspace{0.2em}
    \item Prosedur menekan tiga register tetapi hanya memulihkan dua
          di antaranya. Akibatnya adalah \dots
          \begin{enumerate}\footnotesize
              \item assembler menolak program karena \asm{PUSH} dan
                    \asm{POP} tidak seimbang
              \item \asm{RET} melompat ke alamat yang bukan alamat
                    kembali
              \item register \reg{AX} dipulihkan otomatis oleh DOS
              \item program tetap benar karena \reg{AX} tidak dipakai
                    lagi
          \end{enumerate}
\end{enumerate}
\end{frame}

% =========================================================
% PENUTUP
% =========================================================
\begin{frame}{Penutup}
\small
\begin{block}{Sebelum pertemuan berikutnya}
\begin{itemize}
    \item Selesaikan laporan praktikum: tabel \eng{trace} ketiga
          kasus bug, salinan keluaran \code{R}, \code{U}, \code{D},
          dan \code{T}, serta berkas hasil perbaikan.
    \item Ulangi bagian yang belum dikuasai dengan menelusuri
          \file{bug-loop.asm} dan \file{bug-stack.asm} satu instruksi
          demi satu, lalu bandingkan keluaran \code{T} dan \code{P}.
    \item Baca \textbf{Bab XVI} --- Analisis Kode Mesin, Disassembly,
          dan Optimasi.
\end{itemize}
\end{block}

\vspace{0.3em}
{\footnotesize\textbf{Bacaan lanjutan:} Bab XV bagian Perbedaan
\code{T} dan \code{P}, Membaca Bingkai Stack, dan Metodologi
\eng{Debugging} yang Sistematis; pembahasan register pada Bab VII;
pembahasan \eng{flag} pada Bab IX; pembahasan prosedur dan
\eng{stack} pada Bab XIII.\par}

\vspace{0.35em}
\begin{center}
{\color{obeBlue}\sffamily\large\bfseries Terima kasih\par}
\vspace{0.12em}
{\color{obeGray}\sffamily\small \NamaMK\ (\KodeMK) \textbullet\
\SemesterTeks\par}
\end{center}
\end{frame}

\end{document}
